\section{Implementation}	\label{a:ab_rts_en}

This chapter covers the implementation of BO with a Cheetah-based physics-informed prior for the transverse beam tuning at the ARES Experimental Area. 
Instead of being one large script, the implementation code is systematically divided into Python modules, each having their distinct functions:
\begin{itemize}
    \item \textbf{\texttt{bo\_cheetah\_prior\_ares.py}} includes simulated ARES environment \texttt{ares\_problem()} and the GPyTorch compatible prior mean class \texttt{AresPriorMean} with its trainable  misalignment parameters. This is where most of the physics lives.
    \item \textbf{\texttt{eval\_ares.py}} constructs the optimisers (Standard BO, BO with Cheetah prior and Nelder-Mead) along with the experimental scenarios (matched, mismatched and matched\_prior\_newtask) and loops through the optimisation trials, collecing data and saving raw results. For our evaluation, we run 10 separate trials for each of the four methods (NM, standard BO, BO with mismatched prior and BO with matched\_prior\_newtask) and save the results to CSV files for later aggregated analysis as next step.
    \item \textbf{\texttt{eval\_ares\_metrics.py}} computes the evaluation metrics and generates a final CSV table with all the aggregated metrics to compare the techniques. 
    \item \textbf{\texttt{plot\_ares\_results.py}} generates plots from the results. 
\end{itemize}

The plotting and metrics modules are then connected by the fifth script \textbf{\texttt{run\_evaluation.py}} to run altogether. Table \ref{table:libraries} shows the summary of primary libraries used. 

\begin{table} [htbp]
    \centering
    \small
    \begin{tabular}{p{2.5cm} p{5.5cm} p{1.5cm}}
        \toprule
        \textbf{Libraries} & \textbf{Purpose} & \textbf{Version} \\
        \midrule
        \texttt{Cheetah} & PyTorch based differentiable beam tracking simulator. & 0.8.0 \\
        \texttt{GPyTorch} & Gaussian process framework, base class for prior mean module. & 1.14.2 \\ 
        \texttt{Xopt} & BO loop, UCB acquisition Nelder-Mead baseline. & 2.6.7 \\ 
        \texttt{BoTorch} & BO primitives used internally by Xopt. & 0.16.0 \\ 
        \bottomrule
    \end{tabular}
    \caption{Primary libraries/frameworks used.}
    \label{table:libraries}
\end{table}

\subsection{Initial Setup}
Before any optimiser, we need something to optimise against, \texttt{bo\_cheetah\_prior\_ares.py} contains
the method \texttt{ares\_problem()} which fulfills this purpose. A beam is propagated across a Cheetah model of the ARES Experimental 
Area using a set of magnet parameters, and the beam quality is then returned. In practice, every call to this function 
mimics what would happen if the user adjusted the magnet settings on the real machine and then looked at the diagnostic screen. \

In this project, Cheetah's \texttt{ParameterBeam} is used to describe beam entering the ARES Experimental Area. It is used 
simply because it is faster, which enables to run the experiments more quickly and iterate on the implementation with shorter turnaround times.
 Also, note that Cheetah's \texttt{ParticleBeam} could have also been just as well but with more computational cost. Since speed is not the most critical 
 concern in this project, there is nothing holding us back from using either representation for this application. The ARES experimental area only 
 involves linear optics, so both approaches get equivalent results. These values match the typical ARES operating conditions.

A JSON file describing the full ARES accelerator is used to load the lattice itself. The code extracts the sub-segment between AREASOLA1 (marker) and AREABSCR1 (screen) 
as only ARES Experimental Area is required. Listing \ref{code:ares_loading} demonstrates the loading of the lattice and the assignment of the five magnet parameters prior to tracking.


\lstset{language=Python,
    basicstyle=\ttfamily\small,      
    keywordstyle=\color{blue},        
    commentstyle=\color{teal},       
    stringstyle=\color{purple},       
    showstringspaces=false,          
    breaklines=true,                 
    numbers=left,                     
    numberstyle=\tiny\color{gray},    
    frame=single,                      
    backgroundcolor=\color{lightgray!20},
    xleftmargin=2em,                    
    framexleftmargin=2em                 
}

\begin{lstlisting}[caption={ARES lattice loading and magnet parameters configuration.}, label={code:ares_loading}]

# Load ARES lattice
ares_segment = cheetah.Segment.from_lattice_json("ARESlatticeStage3v1_9.json")
ares_ea = ares_segment.subcell("AREASOLA1", "AREABSCR1")

# Set magnet parameters
ares_ea.AREAMQZM1.k1 = torch.tensor(input_param["q1"])
ares_ea.AREAMQZM2.k1 = torch.tensor(input_param["q2"])
ares_ea.AREAMCVM1.angle = torch.tensor(input_param["cv"])
ares_ea.AREAMQZM3.k1 = torch.tensor(input_param["q3"])
ares_ea.AREAMCHM1.angle = torch.tensor(input_param["ch"])

if misalignment_config is not None: 
    for magnet_name , (dx, dy) in misalignment_config.items():
        magnet = getattr(ares_ea, magnet_name)
        magnet.misalignment = torch.tensor([dx, dy], dtype=torch.float32
        )
else:
        ares_ea.AREAMQZM1.misalignment = torch.tensor([0.0, 0.0])
        ares_ea.AREAMQZM2.misalignment = torch.tensor([0.0, 0.0])
        ares_ea.AREAMQZM3.misalignment = torch.tensor([0.0, 0.0])

out_beam = ares_ea(incoming_beam)
\end{lstlisting}

For the three quadrupole strengths ($k^{Q1}$, $k^{Q2}$, $k^{Q3}$) and corrector angles ($\theta^{CV}$ and $\theta^{CH}$), we have reduced the range from physical magnet's 
\SI{\pm72}{m^{-2}} to our simulation's \SI{\pm30}{m^{-2}}, and from \SI{\pm6.2}{mrad} to \SI{\pm6.0}{mrad} respectively because the optimisation consistently converges well within the narrower range
and exploring beyond it offers no practical benefit to the tuning process, since during practical operation, these ranges are usually never exceeded. The code iterates through the dictionary and applies each ($\Delta_x, \Delta_y$) pair to the corresponding 
quadrupoles when a misaligned configuration is provided. This enables us to model the realistic scenario in which magnets are not perfectly 
centred on the beam. To let the optimiser know how good or bad the result is, we need a single number. Following the previous FODO demonstration \cite{kaiser2024bridging}, for ARES Experimental Area, we use the mean absolute error as objective function defined in Equation \ref{eq_ares_ob}, including each of the four beams observable, with $\mu_x$ and $\mu_y$ being the beam positions (ideally, zero means beam is centred) and $\sigma_x$ and $\sigma_y$ being the beam sizes (the smaller the better), simultaneously centering and focusing the beam. The function also returns the individual beam parameters in addition to MAE so that they can be logged for later analysis.

\subsection{Physics-Informed Prior Mean Module}
The heart of this project work is \texttt{AresPriorMean}, defined in \texttt{bo\_cheetah\_prior\_ares.py}. 
Its role is to turn the Cheetah simulation into a format that GPyTorch can use as the mean function of GP. 
Instead of using GP prior (typically with zero mean assumption), it begins with the best estimates from the physics model and 
just needs to learn the residual, which is the difference between what the simulator predicts and what is actually observed. \

\texttt{AresPriorMean} class specifically inherits from \texttt{gpytorch.means.Mean} and implements a \texttt{forward()} method. 
Everytime GPyTorch needs the prior mean at a set of input points, it calls this method. The remaining GP components, such as kernel evaluation, 
posterior conditioning and acquisition function optimisation, perform exactly like they would with any other mean function. 
Another important aspect is that Xopt builds the input tensor for the GP by sorting variables in alphabetical order, which means that since our 
magnet variables are named $q1$, $q2$, $cv$, $q3$ and $ch$, column 0 of the tensor would  always be $ch$, not $q1$ as one might expect from the beamline order. 
Listing \ref{code:xopt_indexing} shows how the class defines explicit index constants in order to avoid silently feeding the wrong values to incorrect magnets.

\begin{lstlisting}[caption={Explicit indexing to ensure correct mapping in Xopt input tensor.}, label={code:xopt_indexing}]

# Load ARES lattice
class AresPriorMean(Mean):
    # Considering alphabetical order
    VAR_ORDER = ['ch', 'cv', 'q1', 'q2', 'q3']  
    IDX_CH = 0
    IDX_CV = 1
    IDX_Q1 = 2
    IDX_Q2 = 3
    IDX_Q3 = 4
\end{lstlisting}

As discussed earlier, in real accelerators, the quadrupole magnets are never perfectly aligned. A magnet that is even with a fraction of a millimeter, 
deflects the beam in a way that would not be predicted by a perfectly aligned simulation, partially acting as a steering factor.  
In order to handle this, the prior module treats the quadrupole misalignments as learnable parameters ($\Delta_{xQ1}$, $\Delta_{yQ1}$, $\Delta_{xQ2}$, $\Delta_{yQ2}$, $\Delta_{xQ3}$, $\Delta_{yQ3}$).
Each of the three quadrupoles have their respective horizontal and vertical misalignment values, making a total of six values. In order to do online system identification, GPyTorch’s marginal likelihood maximization 
adjusts these values in addition to the standard GP hyperparameters as the optimiser collects data. Three coordinated layers are used to register each misalignment parameter, these three 
layers are summarized in Table \ref{three_layers}. On top of this, $@property$ accessors offer a convenient interface, reading a misalignment property applies the forward transform (raw to constrained) and when 
it is written, the inverse transform is applied. This three layer pattern is repeated identically for each of the six parameters.  The constructor is fairly lengthy because of this six times repeated 
pattern, although a loop based factory would have been cleaner, but we go with the explicit version since it has the advantage of being easy to read and to debug. 

\begin{table} [htbp]
    \centering
    \small
    \begin{tabular}{p{4.0cm} p{4.0cm} p{7.0cm}}
        \toprule
        \textbf{Component} & \textbf{GPyTorch API} & \textbf{What it does} \\
        \midrule
        Raw Parameter & \texttt{register\_parameter()} & A raw unconstrained \texttt{nn.Parameter} is created at first, which is what PyTorch’s optimiser updates directly.\\
        Interval Constraint & \texttt{register\_constraint()} & Secondly, an interval constraint is used that maps this raw value through a sigmoid based transformation into the physically meaningful range of \SI{\pm0.5}{mm}, this serves as the bound guaranteeing that the learnt offset can never be greater than what is practical for quadrupole misalignments. \\ 
        Smoothed box prior & \texttt{register\_prior()} & Third, a \texttt{SmoothBoxPrior} is registered across the confined space, adding a soft regularization on top of the hard constraint. This helps the optimiser avoid getting stuck at the boundaries by gently penalizing values near the edges while assigning a fairly uniform probability throughout the interval. \\ 
        \bottomrule
    \end{tabular}
    \caption{Three layered pattern to register misalignments as trainable parameters.}
    \label{three_layers}
\end{table}

When \texttt{forward()} method is called, it recieves a tensor X whose last dimension contains 5 columns (the five magnets, in alphabetical order), the steps include:
\begin{enumerate}
    \item Using the index constants from Listing \ref{code:xopt_indexing}, slice out each magnet parameter. 
    \item Assign them to corresponding components of Cheetah lattice.
    \item Apply the current misalignment values to the three quadrupoles by stacking them.
    \item Run the beam through the lattice
    \item Return the MAE (Equation \ref{eq_ares_ob}).
\end{enumerate}
Since all of these operations (the beam tracking, the absolute values, and the summation) are native PyTorch operations, the full forward pass is differentiable. 
Gradients could flow backward through the simulator and into the six misalignment parameters, which is what allows the joint optimisation of both misalignment 
values as well as usual GP hyperparameters.

\subsection{Optimisation Runner}

After setting up the environment and the prior setup, the remaining part is the runner that actually executes the experiments. This is handled by an \texttt{eval\_ares.py} file 
that takes arguments for the experimental scenarios (matched, mismatched, matched\_prior\_newtask), the number of trials, the evaluation budget and the type of optimiser (BO, BO\_prior, NM) 
where BO is the standard Bayesian Optimisation, BO\_prior is Bayesian Optimisation with physics prior and NM is Nelder-Mead. Xopt defines the optimisation problem declaratively using a 
YAML based specification known as VOCS (Variable and Objective Configuration Specification). Listing \ref{code:VOCS_configuration}  shows the five magnet's ranges and a single objective.


\begin{lstlisting}[caption={VOCS configuration.}, label={code:VOCS_configuration}]
vocs_config = """
        variables:
            q1: [-30, 30]
            q2: [-30, 30]
            cv: [-0.006, 0.006]
            q3: [-30, 30]
            ch: [-0.006, 0.006]
        objectives:
            mae: minimize
    """
    vocs = VOCS.from_yaml(vocs_config)
\end{lstlisting}

To test the prior under various conditions, the runner supports three scenarios given below, 
they are also summarised in Table \ref{table:scenarios_summary}:

\begin{enumerate}
    \item \textbf{Matched Scenario:} This is the most simple and ideal scenario where the beam and misalignments in the environment are exactly what the prior expects, which are the default values (including all misalignment values to be zero). Thus, physics model perfectly matches the reality. This scenario exists mainly to verify that the implementation works correctly and to see how much the prior helps when it is perfectly accurate since there is no model reality gap at all. In practice, this situation never happens on a real machine. 
    \item \textbf{Mismatched Scenario:} In this scenario, a gap is introduced on purpose to create mismatch between the prior and the reality, a different beam ($\sigma_y$ changed from \SI{200}{\mu}m to \SI{100}{\mu}m), and non zero misalignment values over all three quadrupoles are used, shown in Table \ref{table:misalignment_groundtruth}. The prior, however, starts from the default beam and zero misalignments as it does not know about any of these changes. We tell Xopt to include the six misalignment parameters in the GP’s marginal likelihood optimisation by setting \texttt{trainable\_mean\_keys=["mae"]} in the \texttt{StandardModelConstructor} so they can be learned from data as it is received. This is the scenario that tests whether online systems identification actually works. 
     \begin{table} [htbp]
    \centering
    \small
    \begin{tabular}{p{4.0cm} p{2.0cm} p{2.0cm}}
        \toprule
        \textbf{Magnet} & \textbf{$\Delta_x$ (mm)} & \textbf{$\Delta_y$ (mm)} \\
        \midrule
        AREASMQZM1 & 0.000 & +0.200 \\
        AREASMQZM2 & +0.100 & -0.300 \\ 
        AREASMQZM3 & -0.100 & +0.150 \\ 
        \bottomrule
    \end{tabular}
    \caption{Ground truth quadrupole misalignments.}
    \label{table:misalignment_groundtruth}
\end{table}


    \item \textbf{Matched\_prior\_newtask Scenario:} The simulated accelerator setup is the same as in the mismatched scenario (modified beam and non-zero misalignments), but this time the prior is initialised with the correct values which are the right beam parameters and the true misalignment values. So this way, we are handing the prior a perfectly adjusted physics model and asking how fast can the optimiser converge when it already knows the true system’s state, compared to having to learn it during the process.
\end{enumerate}

For the evaluation in this project, the matched\_prior\_newtask serves as our actual matched case, since it uses the realistic 
accelerator setup having modified beam and misalignments rather than the ideal zero misalignment defaults. We compare four methods, 
Nelder-Mead, Standard BO (zero mean), BO with Cheetah prior (mismatched scenario) and BO with Cheetah prior (matched\_prior\_newtask scenario). 
All four approaches face the same simulated accelerator so that the comparison is fair,  what differs between them is only the optimisation strategy, 
and in case of prior informed methods, how much the prior knowledge is available from the start for each method. 



Listing \ref{code:mismatched_setup} shows how the most complex one (BO with mismatched prior) is set up.
the others (NM, standard BO and BO matched prior) are the simpler subsets of the same pattern. 

\begin{lstlisting}[caption={BO with Cheetah prior(mismatched setup).}, label={code:mismatched_setup}]
# Create the prior mean module with default
prior_mean_module = bo_cheetah_prior.AresPriorMean(
                    incoming_beam=incoming_beam
                )
                # Initialise misaligments to zero, will be learned
                prior_mean_module.q1_misalign_x = 0.0
                prior_mean_module.q1_misalign_y = 0.0
                prior_mean_module.q2_misalign_x = 0.0
                prior_mean_module.q2_misalign_y = 0.0
                prior_mean_module.q3_misalign_x = 0.0
                prior_mean_module.q3_misalign_y = 0.0 

                #Inject into GP constructor with trainable_mean_keys
                gp_constructor = StandardModelConstructor(
                    mean_modules={"mae": prior_mean_module},
                    trainable_mean_keys=["mae"], 
                )    
                
                # Create UCB generator with beta 2.0
                generator = UpperConfidenceBoundGenerator(
                beta=2.0, vocs=vocs, gp_constructor=gp_constructor
            )    
\end{lstlisting}


\begin{table} [htbp]
    \centering
    \small
    \begin{tabular}{p{2.5cm} p{2.0cm} p{2.0cm} p{2.0cm} p{2.0cm}}
        \toprule
        \textbf{Scenario} & \textbf{Modified Beam} & \textbf{True Misalignments} & \textbf{Prior Misalignments} & \textbf{Trainable $\phi$} \\
        \midrule
        Matched & --- & --- & --- & --- \\
        Mismatched & $\checkmark$ & $\checkmark$ & 0 & $\checkmark$ \\ 
        Matched prior (new task) & $\checkmark$ & $\checkmark$ & = true & ---  \\ 
        \bottomrule
    \end{tabular}
    \caption{Three experimental scenario summary.}
    \label{table:scenarios_summary}
\end{table}

In order to have a fair comparison, every trial starts from the same initial points, which are

\begin{equation}
    x_{0} = (k^{Q1}=10.0, k^{Q2}=-10.0, \theta^{CV} = 0.0, k^{Q3} = 10.0, \theta^{CH} = 0.0).
\end{equation}

After this initial evaluation, the runner loops through the \texttt{xopt.step()} for a fixed number of iterations. 
The default in the code is 50 steps, but for the evaluation presented in this report, we used 100 steps to give each 
method more room to converge and to better observe the later stage behaviors of the optimisers. Each step involves 
fitting the GP to all the data collected so far, selecting next point by optimising the UCB acquisition function 
and testing that point on the simulated accelerator. In the case of Nelder-Mead, the same \texttt{step()} interface 
updates the simplex instead. For mismatched scenario when the prior’s misalignment parameters are trainable, fitting 
the GP also involves optimising over \textbf{$\phi$}. As a result, each new data point gives the marginal likelihood 
optimiser more information about where the magnets are actually positioned and the prior gradually gets more accurate 
which is the online system identification already discussed above. After each trial completes, the script stores the 
full MAE history, the overall best MAE, beam parameters and the learned misalignment values (where applicable). With a fixed random seed, 
10 separate trials are conducted per configuration and the results are saved to CSV file for further analysis during evaluation metrics computation. 


\subsection{Evaluation Metrics}
The evaluation code is located in \texttt{eval\_ares\_metrics.py}. 
For each metric, $i \in \{1,2,\ldots,10\}$, where $i$ denotes the trial index, $t \in \{0,1,\ldots,T\}$, where $t$ is the step (iteration) index and $T$ is the last step which is 100 in our experiments.
\begin{enumerate}
    \item \textbf{Best MAE:} is the best (minimum) mean absolute error (MAE) reached at any step within each trial, defined as 
    \begin{equation}
        \text{Best MAE}^{(i)} = \min_{t \in \{0, 1, \ldots, T\}} \text{MAE}_t^{(i)}.
    \end{equation}
    \item \textbf{Final MAE:} is the MAE at the last step of each trial ($\text{MAE}_T^{(i)}$).
    \item \textbf{RMSE:} is the root mean square error across all iterations within each trial, capturing how quickly the optimiser improves and not just the final value, defined as
    \begin{equation}
        \text{RMSE}^{(i)} = \sqrt{\frac{1}{4(T+1)} \sum_{t=0}^{T} \left( (\mu_x^{(i,t)})^2 + (\mu_y^{(i,t)})^2 + (\sigma_x^{(i,t)})^2 + (\sigma_y^{(i,t)})^2 \right)}.
    \end{equation}
    \item \textbf{Improvement Best percentage(\%):} is the improvement measured from initial MAE to best MAE from each trial $i$, computing as $(((\text{MAE}_0 - \text{MAE}_{\text{best}})/\text{MAE}_0) \times 100)$.
    \item \textbf{Improvement Final percentage(\%):} is the improvement measured from initial MAE to Final (last step) MAE of each trial $i$, computing as $(((\text{MAE}_0 - \text{MAE}_{\text{final}})/\text{MAE}_0) \times 100)$.

    \item \textbf{Steps to target:} computes how many evaluations does it take to drop below the threshold of \SI{40}{\mu}m of MAE for the first time, this shows how quickly an optimiser gets the beam into an acceptable state.
    \item \textbf{Success rate:} is the percentage of trials that reach the 40 $\mu$m target within the budget.
\end{enumerate}

After computing all the metrics for each trial $i$, we then compute and report the mean, median and standard deviation of those metrics across the 10 trials.
Metric definition and evaluation are adapted from~\cite{kaiser2024bridging, kaiser2024rl}.